\documentclass[10pt,a4paper]{amsart}
\usepackage[margin=19mm]{geometry}
\usepackage{amsmath,amssymb,mathtools}
\usepackage[scr=rsfs,cal=euler]{mathalfa}
\usepackage{xfrac}
\usepackage[unicode,hidelinks,bookmarks=false]{hyperref}
\newcommand{\ie}{i.e.\ }
\newcommand{\cf}{cf.\ }
\newcommand{\resp}{respectively\ }
\newcommand{\Subsection}{\subsection}
\providecommand{\nobreakdash}{\nobreak\hbox{-}}
\setlength{\emergencystretch}{2em}
\multlinegap0pt
\usepackage{mathtools}
\usepackage{amssymb}
\usepackage{enumerate}
\newtheorem{lemma}{Lemma}
\usepackage{cleveref}
\crefname{lemma}{Lemma}{Lemmas}
   \newcommand{\bl}{\otimes_\zl\fl}
   \newcommand{\df}{\dim_\fl}
 \newcommand{\fl}{{\overline{\mathbb{F}}_\ell}}
   \newcommand{\fo}{\mathfrak{o}}
   \newcommand{\ho}{\mathrm{Hom}}
   \newcommand{\id}{\mathrm{Ind}}
   \newcommand{\im}{\mathrm{Im}}
   \newcommand{\ip}{\mathrm{Ind}_P^G}
   \newcommand{\ipp}{\id_{P'}}
   \newcommand{\oo}{\mathbb{O}}
   \newcommand{\op}{\overline{P}}
   \newcommand{\pad}{\mathfrak{d}}
 \newcommand{\ql}{{\overline{\mathbb{Q}}_\ell}}
   \newcommand{\ra}{\rightarrow}
  \newcommand{\rep}{\mathrm{Rep}}
   \newcommand{\res}{\mathrm{res}}
   \newcommand{\tpi}{\Tilde{\pi}}
   \newcommand{\xc}{\mathfrak{X}}
  \newcommand{\zl}{{\overline{\mathbb{Z}}_\ell}}
\title{Lifting banal representations of classical groups}
\author{Johannes Droschl}
\date{Source: arXiv:2604.07516v1 (CC BY 4.0)}
\begin{document}
\maketitle

\noindent\emph{Source and licence.} Lifting banal representations of classical groups. Version arXiv:2604.07516v1 (CC BY 4.0), distributed by arXiv under the Creative Commons Attribution 4.0 International licence, http://creativecommons.org/licenses/by/4.0/, as recorded in the arXiv licence metadata for this exact version and shown on the abstract page https://arxiv.org/abs/2604.07516v1. Full licence notice: \texttt{licenses/arxiv-2604.07516-CC-BY-4.0.txt}.\par\smallskip

\noindent\textit{Editorial scope.} Counted unit: the article's lemma L:ratJ (source0.tex lines 454--462). The article's complete original proof of it is delivered with it (source0.tex lines 463--477, one \texttt{proof} environment of 15 lines). No mathematics is written, repaired or re-derived: the statement and the proof are the article's own lines, in the article's own notation, and no retained line is substituted or re-flowed.  Retained uncounted as the source-local context the proof needs: lines 372--379; lines 381--384.  Nothing was omitted from any retained span, so the package carries no omission record.  No retained line contains a citation, so the wrapper carries no bibliography and no bibliography entry.  No retained line contains a figure environment, an image-inclusion command or drawing code, so no image is delivered and the package declares no asset dependency.  Editorial formatting only, in addition: an AMS \texttt{amsart} preamble that restates the article's own declarations and macros used by the retained lines, namely the environments \texttt{lemma} on separate counters, together with the standard packages those lines need.  The retained environments print in this document as Lemma 1 in order of appearance, whereas the article prints the same items with the numbering of its own full document.  The complete native source of the article as distributed in the arXiv e-print for this exact version is bundled unchanged as \texttt{sources/198094/source0.tex}.
\begin{lemma}\label{L:pder}
    For $\pi$ and $P$ as above we have $\pad(\pi,\op)=\pad(\pi,P)\ge 0$. Moreover, the following are equivalent.
    \begin{enumerate}
        \item $J_{P,\op}(\pi)$ is an isomorphism.
        \item $J_{\op,P}(\pi)$ is an isomorphism.
        \item $\pad(\pi,P)=0$.
    \end{enumerate}
\end{lemma}

\begin{lemma}\label{L:noresidue}
    Let $\pi$ and $P,P'$ be as above. Then $\Lambda(\pi,P,P')=0$ if and only if the map
    $r_{P'}(\ip(\pi))\ra \pi$ obtained by Frobenius reciprocity does not vanish on $r_{P'}(F(\oo'))(\pi)\subseteq r_{P'}(\ip(\pi))$ if $\oo'\subsetneq\overline{\oo}$. 
\end{lemma}

\begin{lemma}\label{L:ratJ}
    Let $\tpi\in\rep_\ql(M)$, $\fo$ be an integral structure of $\tpi$, and write $\pi\coloneqq\fo\bl$. Assume that the following is satisfied.
    \begin{enumerate}
        \item $\df\ho(\ip(\pi),\id_{\op}^G(\pi))=1$.
        \item $\Lambda(\tpi,P,\op)=\Lambda(\pi,P,\op)=0$.
    \end{enumerate}
    Then \[J_{P,\overline{P}}(\tpi)(\ip(\fo))\subseteq \id_{\op}^G\fo).\]
    Moreover, if for all $1\neq a\in\zl^\times$ with $a=1\mod\ell$, we have that \[\Lambda(\tpi\otimes\nu_\alpha,P,\op)=\alpha(\tpi\otimes\nu_\alpha,P)=0,\] then \[J_{P,\overline{P}}(\tpi)(\ip(\fo))\bl=\im(J_{P,\op}(\pi)).\]
\end{lemma}

\begin{proof}
We start by proving the first inclusion. Indeed, from the Bauer-Nesbitt principle of Vigneras, it follows that there exists $k\in\mathbb{N}$ such that the map $I\coloneqq \ell^kJ_{P,\overline{P}}(\tpi)$ satisfies the inclusion, and we fix a minimal such $k$. We want to show that $k=0$. We thus can reduce $I$ mod $\ell$ to obtain a non-zero map $I\bl\in \ho(\ip(\pi),\ipp(\pi))$, which implies that up to a scalar $I$ equals to $J_{P,\overline{P}}(\pi)$. Note that by \Cref{L:noresidue} the latter does not vanish for some $f\in J_{P,\overline{P}}(\pi)$ with support in $P\op$. Similarly, by \Cref{L:noresidue}, it follows that the same statement holds for $J_{P,\overline{P}}(\fo)$. Now if $k>0$, it would follow that $I$ vanishes on all $f\in J_{P,\overline{P}}(\tpi)$ with support in $P\op$, a contradiction. Thus $k=0$ and the claim follows.

For the second claim we use  the rationality of intertwining operators. It suffices to prove that for $f\in \ip(\fo)$ with $f\mod\ell\neq 0$, $J_{P,\overline{P}}(\tpi)(f)=0\mod \ell$ implies $J_{P,\op}(\tpi)(f)=0$. 
Let $K'$ be a sufficiently small open compact subgroup of $G$ fixing $f$.
We therefore find a basis $\{h_1,\ldots,h_k\}$ of $\id_{P\cap K}^K(\tpi)^{K'}$  and $\{f_1,\ldots,f_k\}$ of $\id_{\op\cap K}^K(\tpi)^{K'}$
such that there exists a $k\times k$ matrix $P=(P_{i,j})$ with entries in $R[T,T^{-1}]$ such that \[\res_K(J_{P,\op}(\tpi\otimes\nu_a)(\res_K^{-1}(f_j)))=\sum_iP_{i,j}(a)f_j.\]
We can assume without loss of generality that the $h_i$'s and $f_i$'s have values in $\fo$ and hence for every $\zl$-valued unramified character $\nu_a$ we have that $P_{i,j}(\nu_a)\in\zl$.
We let \[\nu\colon \xc_\ql[M]\ra \ql[T,T^{-1}]\] be the ring map used in the definition of intertwining operators.
 By assumption, the matrix $\nu(P)(1)\mod \ell$ has non-maximal rank and hence $\det(\nu(P)(1))=0\mod\ell$.
 Let $\alpha_1,\ldots,\alpha_m$ be the zeros of $\det(\nu(P))$ over $\ql$. Since $\det(\nu(P)(1))\in \zl$ and equals to $0$ mod $\ell$, at least one of the $\alpha_i$  must be integral and reduce to $1\mod \ell$.
 Therefore, there exists an integral character $\nu_{\alpha_i} $, which lifts the trivial character over $\fl$ to $\zl$, such that $\det(P(\nu_{\alpha_i}))=0$. But then this implies that $J_{P,\op}(\tpi\otimes\nu_{\alpha_i})$ is not an
  isomorphism, which in turn implies via \Cref{L:pder} and the assumption of the lemma that $\alpha_i=1$.
  Hence $J_{P,\op}(\tpi)(f)=0$ as desired. 
\end{proof}

\end{document}
